\section{The Russian Idea}

\begin{quotex}
For as lightning cometh out of the east, and appeareth even into the west: so shall the coming of the Son of man be. 
\flright{\textsc{Matthew 24:27}}
\end{quotex}


Over the past two hundred years, an important spiritual and philosophical movement has been coming out of Russia. This, until recently, has been little known and appreciated; fortunately, the past few decades has seen new translations of many of the thinkers listed below. According to the Lindisfarne Press\footnote{\url{http://www.steinerbooks.org/p.php?id=12}}, the publishers of the \textit{Library of Russian Philosophy}, Russian philosophy is marked by certain characteristic features.

\begin{itemize}


\item \textbf{Epistemological realism:} Knowledge consists of the idea of the thing, and these ideas --- in the Platonic sense --- are real.

\item \textbf{Integral knowledge:} Knowledge is an all-embracing unity that includes sensuous, intellectual, and mystical intuition. Here we recognize the degrees of knowledge described by Rene Guenon, in recapitulation of Traditional Wisdom.

\item \textbf{Integral personality:} The complete man is sensuous, rational, and mystical.

\item \textbf{Resurrection:} The flesh and the world are to be transformed.
\end{itemize}


I date this from the time of the Joseph de Maistre's\footnote{\url{https://gornahoor.net/library/EvolaOnDeMaistre.pdf}} period as ambassador to Russia. In addition to his vocation, Maistre was a Martinist and Catholic philosopher. He represents the opening of Russia to Western ideas, and around the same time, several works important to Martinism and Hermeticism were translated into Russian. These include: the \textit{Imitation of Christ} by Thomas a Kempis, as well as works by Louis-Claude de Saint-Martin and Jacob Boehme.

This led to two independent streams.

\begin{itemize}


\item Occult, as represented by

  \begin{itemize}

  
\item Helena Blavatsky
  
\item Gurdjieff, Ouspensky
  \end{itemize}


\item Sophiology, as represented by

  \begin{itemize}

  
\item Vladimir Solovyov
  
\item Sergius Bulgakov
  
\item Pavel Florensky
  
\item Nikolai Berdyaev
  
\item Vyacheslav Ivanov
  \end{itemize}

\end{itemize}


This spiritual outpouring can related to the three Romes.

\begin{itemize}


\item The Third Rome. Theophan the Recluse ties together the two Romes: one, by the Russian edition of the Philokalia, and the other in his version of \textbf{Lorenzo Scupoli}`s \emph{Spiritual Combat}, translated into English as \emph{Unseen Warfare}.

\item Boris Mouravieff represents the second Rome, the Byzantine Empire. Following Theophan and the occultism of Gurdjieff and Ouspensky, he reworks the Orthodox spiritual tradition.

\item Valentin Tomberg is the first Rome. He begins with the occult stream originating with Blavatsky and then expanded by Rudolf Steiner. He is also heir to the Martinist current that had infiltrated the Russian consciousness and then integrates it with the Sophiology of Solovyov and his followers. He eventually joined the Roman Church.
\end{itemize}


I include Rudolf Steiner, although he was not a Russian However, his wife was Russian, and he predicted that the next cultural epoch is to emerge from Russia.

The culmination of these streams is represented by two important works of Christian Hermeticism:

``Meditations on the Tarot'' by Valentin Tomberg, and \emph{Gnosis}, by Boris Mouravieff. While quite different in tone and subject matter, they are compatible on a deeper level. I can point out some interesting correspondences.

\begin{itemize}


\item Both books written by Russians, who were almost exact contemporaries: Valentin Tomberg (1900 -- 1973) and Boris Mouravieff (1890 -- 1966).

\item They each fled the Russian revolution in 1920.

\item They each claimed to be recovering the esoteric teachings of the Church, one from the Roman church, the other from the Eastern church.

\item Their main works were written in the 1960s in French.

\item They were made available in English within a few years of each other.
\end{itemize}


The first two volumes of \emph{Gnosis} (exoteric and mesoteric cycles) each contains 21 chapters, indicating incompleteness. The third (esoteric cycle) has 22 chapters. It concludes with a description of the Knight, his work in the world, and the creation of an elite to begin the process of transformation.

\emph{Meditations on the Tarot} contains 22 chapters, one for each of the major arcana. The final arcana is the ``world'', and it ends with the challenge to bring Hermeticism into the world by meditating on the meanings of the minor arcana.


\flrightit{Posted on 2010-03-27 by Cologero}

